\section*{September 4, 2026: Baseline demographics and transaction weighting}
\textit{Agent-drafted research entry; interpretations remain open for discussion.}

Jacob requested neighborhood demographics for the comparison of 30 closed
and 49 still-open schools on the February 2013 candidate list. All 10,921
eligible transactions match
exactly one 2012 Census tract and one block group using their home locations.
The pre-closure sales span 181 distinct tracts and 316 distinct block groups.
The audit attaches 2008--2012 American Community Survey (ACS) five-year
estimates to each home's Census area. These describe neighborhood conditions
before closure and remain fixed across sale years.

\begin{center}
\small
\input{../output/demographic_table.tex}
\end{center}
{\footnotesize Each 2008--2012 eligible transaction receives its tract's ACS
characteristic. The columns average those values over transactions. Income
is a mean of tract medians, not the median income of pooled residents.}

Treated sales occur in neighborhoods with more college graduates, lower
unemployment, and greater population density. Baseline vacancy is nearly
identical across groups. Block groups yield similar directional differences
for education, unemployment, and racial composition. Many differences shrink
when every school site receives equal weight, after averaging its nearby
sales. This supports examining the contribution of sales across sites; it
does not establish parallel trends under either weighting scheme.

Codex recommends equal transaction weights for the baseline. Weighting by
sale value would use the outcome as a weight and amplify expensive sales;
weighting by square footage would change the target toward larger properties.
A fixed pre-2013 sales-count weight for each site is a possible sensitivity,
not an adopted specification. Equal site weighting in this audit first
averages baseline characteristics across each site's pre-closure sales and
then gives every site equal weight within its treatment group.

At Jacob's request, the balance plot adds 95\% percentile bootstrap intervals
using 1,999 draws of school sites within treatment status. Each draw retains
the sampled site's sales and recomputes the standardized difference. These
intervals hold Census estimates fixed and assume independence across sites;
they exclude ACS survey error and dependence across nearby sites. They are
not tests of causal identification. The median reported income margin of
error across observed geographies is \$8,385 for tracts and \$13,592.50 for
block groups, supporting tracts as a less noisy first adjustment.

The demographic variables are available for future baseline-by-time controls;
they were not added to the regressions in this audit. Baseline vacancy and
poverty describe initial conditions. Their post-closure changes could be
outcomes or mechanisms and should not automatically become controls.


\clearpage
\begin{center}
\includegraphics[width=\textwidth]{../input/demographic_balance.png}
\end{center}
{\footnotesize Standardized differences divide the treated-minus-control mean
by the root mean of the two group variances. Variances are computed across
transactions for equal transaction weights and across site means for equal
site weights. The denominator is recomputed in each bootstrap draw.}

{\footnotesize Reproduction: run Make in the home-baseline-demographics audit's
code directory, then in the logbook code directory. Inputs are the unchanged
clean sales and exposure files plus the Census acquisition audit's official
ACS bulk sequences and 2012 TIGER files, downloaded September 4. C17002 supplies
poverty because B17001 is unavailable at block-group level in this vintage.
All ten covariates are observed for every pre-closure sale at both levels.
Source cell margins of error and standard dataset reports are retained.
An authenticated API check matched all 2,628 available tract income
estimate/MOE cells to the bulk files. Revision c200c31
plus working changes; the bootstrap seed is 20260904.}
